\chapter{Land surface: the Noah model}\label{ch:lsm}

Over land, the surface fluxes of heat and moisture depend on the state of the soil, the vegetation and the snow. The
case uses the Noah land-surface model \citep{chen2001,ek2003} (\code{sf_surface_physics = 2}). Noah predicts:
\begin{itemize}
  \item the soil temperature and moisture in four layers (0--10, 10--40, 40--100 and 100--200~cm);
  \item canopy water;
  \item snow water equivalent and snow depth;
  \item the skin temperature.
\end{itemize}
Its physics is in \code{phys/module_sf_noahlsm.F}, called per point by the driver
\src{phys/module_sf_noahdrv.F}{lsm}. Glaciers have a variant (\code{module_sf_noahlsm_glacial_only.F}) and sea ice has
its own model (\code{module_sf_noah_seaice.F}).

\section{The structure of one point}

For each land point and step, the driver gathers the forcing and calls \src{phys/module_sf_noahlsm.F}{SFLX}. The
forcing is:
\begin{itemize}
  \item from the lowest model level: wind, temperature, humidity, pressure;
  \item from radiation: the downward short- and longwave fluxes;
  \item from the microphysics: precipitation;
  \item from the surface layer: the exchange coefficients.
\end{itemize}
SFLX calls a tree of routines:
\begin{center}
\begin{tabular}{ll}
\toprule
Routine & Role \\
\midrule
\code{REDPRM} & parameters for the point's vegetation and soil type (from \code{VEGPARM.TBL}, \code{SOILPARM.TBL}, \code{GENPARM.TBL}) \\
\code{CSNOW}, \code{SNOW_NEW}, \code{SNFRAC}, \code{SNOWZ0}, \code{ALCALC} & snow density and depth, snow cover fraction, roughness and albedo \\
\code{TDFCND} & soil thermal conductivity from moisture and soil type \\
\code{PENMAN} & potential evaporation \\
\code{CANRES} & canopy resistance \\
\code{NOPAC} / \code{SNOPAC} & the energy and water budget without or with snow \\
\code{EVAPO}, \code{DEVAP}, \code{TRANSP} & total evaporation: bare soil, transpiration, canopy \\
\code{SMFLX}, \code{SRT}, \code{SSTEP}, \code{WDFCND} & soil water: infiltration, runoff, diffusion, implicit step \\
\code{SHFLX}, \code{HRT}, \code{HSTEP}, \code{TBND}, \code{TMPAVG} & soil heat: diffusion, implicit step \\
\code{FRH2O}, \code{SNKSRC} & frozen soil water and the latent heat of freezing \\
\code{ROSR12} & the tridiagonal solver used by both implicit steps \\
\bottomrule
\end{tabular}
\end{center}

\section{Energy balance and potential evaporation}

The surface energy balance
\begin{equation}
  (1 - \alpha)\,S\!\downarrow + \epsilon L\!\downarrow - \epsilon\sigma T_s^4 = H + \lambda E + G
\end{equation}
is linearized around the air temperature, in the manner of Penman. \code{PENMAN} computes the potential evaporation:
the evaporation of a saturated surface with the available energy and the surface-layer exchange coefficient. Actual
evaporation is the potential evaporation reduced by the soil moisture and the canopy resistance. The skin temperature
then follows from the balance.

\section{Canopy resistance}

Transpiration is limited by the stomata. Noah uses a Jarvis-type resistance (\src{phys/module_sf_noahlsm.F}{CANRES}):
\begin{equation}
  R_c = \frac{R_{c,\min}}{\mathrm{LAI}\;F_1\,F_2\,F_3\,F_4},
\end{equation}
with four factors between 0 and 1:
\begin{itemize}
  \item solar radiation: $F_1 = (f + R_{c,\min}/R_{c,\max})/(1 + f)$, where $f$ grows with the incoming solar flux
    (\code{RCS});
  \item air temperature: $F_2 = 1 - 0.0016\,(T_{\mathrm{opt}} - T)^2$ (\code{RCT});
  \item vapor deficit: $F_3 = 1/\bigl(1 + h_s\,(q_{\mathrm{sat}} - q)\bigr)$ (\code{RCQ});
  \item soil moisture: $F_4$, a root-weighted sum of $(\theta - \theta_w)/(\theta_{\mathrm{ref}} - \theta_w)$ over the
    root layers (\code{RCSOIL}).
\end{itemize}
Each factor has a small positive floor, so $R_c$ stays finite.

\section{Soil heat and water}

Soil temperature obeys the heat diffusion equation with a conductivity that depends on soil moisture:
\begin{equation}
  C(\theta)\,\pd{T}{t} = \pd{}{z}\Bigl(K_T(\theta)\,\pd{T}{z}\Bigr),
\end{equation}
with the ground heat flux at the top and a fixed deep temperature at 8~m. Soil moisture obeys the diffusive form of
Richards' equation:
\begin{equation}
  \pd{\theta}{t} = \pd{}{z}\Bigl(D(\theta)\,\pd{\theta}{z}\Bigr) + \pd{K(\theta)}{z} + F_\theta,
\end{equation}
with infiltration at the top, gravitational drainage at the bottom and sinks from evaporation and root uptake.

Both equations are discretized on the four layers and stepped implicitly in time. Each gives a tridiagonal system per
point, solved by \code{ROSR12}. When the soil freezes, \code{FRH2O} computes how much of the water stays liquid,
and the latent heat enters the heat equation.

\section{Snow}

With snow on the ground, \code{SNOPAC} replaces \code{NOPAC}:
\begin{itemize}
  \item the snowpack is a single layer that accumulates precipitation (\code{SNOW_NEW}, with a temperature-dependent
    density) and compacts (\code{SNOWPACK});
  \item it melts or sublimates from the surface energy balance;
  \item a snow-cover fraction (\code{SNFRAC}) blends the albedo and fluxes of the snow-covered and bare parts.
\end{itemize}

\begin{algorithm}[ht]
\caption{One Noah point step (\code{SFLX}, simplified).}
\label{alg:noah}
\begin{algorithmic}[1]
\State parameters of the vegetation and soil type (\code{REDPRM})
\State snow: new snow, density, cover fraction, albedo, roughness
\State thermal conductivity of the top layer (\code{TDFCND})
\State potential evaporation (\code{PENMAN}); canopy resistance (\code{CANRES})
\If{no snow}
  \State \code{NOPAC}: evaporation components; soil water implicit step; soil heat implicit step; skin temperature
\Else
  \State \code{SNOPAC}: snow energy balance, melt, compaction; then soil water and heat below the pack
\EndIf
\State runoff, fluxes, diagnostics
\end{algorithmic}
\end{algorithm}

\begin{portnote}
Noah is a per-point model with a deep call tree of about thirty routines. Three properties shape its GPU port.
\begin{itemize}
  \item \textbf{Module state.} The point's indices \code{iloc}, \code{jloc} are module variables, and the land-use
    and soil types are passed around as character strings (\code{LUTYPE}, \code{SLTYPE}). The port makes the indices
    arguments and the types integer codes. These are shared refactors, bit-neutral by construction.
  \item \textbf{Implicitly saved locals.} \code{REAL :: SNUPGRD = 0.02} in \code{SNFRAC} is implicitly \code{SAVE}d in
    Fortran. On a GPU, all threads would share it. It becomes a \code{PARAMETER}, which is the same value.
  \item \textbf{Tables.} The vegetation, soil and general parameter tables (49 arrays and scalars) are uploaded to the
    device once and checked by checksum (T-TAB).
\end{itemize}
Each point then runs \cref{alg:noah} in one GPU thread. A point harness (T-NOAH-PT) compares $10^5$ points from
real states between host and device before the model-level tests (\cref{ch:subsys}).
\end{portnote}
